
\chapter{METHODOLOGY}

\section{System Architecture}
\noindent\hspace{2.2em}The system is designed as a web-based platform built on a multi-layered architecture that separates into three distinct layers: the Presentation Layer, the Application Layer, and the Data Management Layer. This separation ensures modularity,maintainability, and a clear division of responsibilities across the system.
\begin{figure}[H]
	\centering
	\includegraphics[width=\textwidth]{System Architecture.png}
	\caption{System Architecture}
	\label{fig:sysarch}
\end{figure}

\subsection{Frontend Layer}
The Frontend of the system is built using React, a component-based JavaScript library that enables dynamic and responsive user interfaces. Each of the three core functions — Schedule Matcher, Rate-My-Course, and Super Planner — is implemented as an independent module where users can navigate between features.

The interface is designed with a focus on simplicity and visual clarity, which is especially important for the Super Planner's drag-and-drop canvas and the Schedule Matcher's result display. React state management handles real-time UI updates, such as rendering matched free-time slots after scheduled images are processed, or updating course ratings as new submissions come in.

The Authentication Session UI is handled at the presentation layer, where users see and interact with related to their login state, for example: the login page, displaying the logged-in student's name or ID on the taskbar, a logout button, disabling or hiding buttons based on whether the user is logged in, and protected routes redirected.

\subsection{Backend Layer}
The Backend is built using Node.js and Express.js, forming a RESTful API server that handles all communication between the frontend and the database. Express.js routes are organized by function, with separate route groups for schedule matching, course rating, and academic planning. 

Image processing for the Schedule Matcher is handled entirely within this layer using JavaScript-native libraries — Tesseract.js for OCR and OpenCV.js for color segmentation This keeps the entire stack in a single language, simplifies deployment, and reduces points of failure. 

Authentication is also enforced at this layer through Express.js middleware, which validates tokens on all protected API routes and rejects unauthorized requests. This ensures that only EGCI students can access and contribute to the platform.
\begin{itemize}
	\item Verifying credentials (checking the submitted student ID/password against the database)
	\item Issuing tokens (e.g. JWT) after successful login
	\item Protecting API routes via middleware that checks for a valid token on every request
\end{itemize}
As shown in Figure \ref{fig:auth}
\newpage
\begin{figure}[H]
	\centering
	\includegraphics[width=\textwidth]{Authentication.pdf}
	\caption{System Authentication Flow}
	\label{fig:auth}
\end{figure}


\subsection{Data Management}
The system uses MongoDB, a NoSQL document-oriented database, as its primary data store practice. MongoDB’s flexible schema design suits well with our web platform since all the three functions store structurally different data. 

Course reviews contain nested rating objects and optional comment fields, while academic plans store arrays of semester slots each containing course objects with prerequisite relationships.

Collections are organized around the three core functions: a schedules collection for temporarily storing extracted schedule data during matching session, a courses collection for course metadata and aggregated ratings, a reviews collection for individual student submissions, and a plans collection for saved Super Planner layouts.

\subsection{Overall Data Flow}
When a user interacts with the React frontend, requests are sent to the Express.js API server. Depending on the request type, the backend either queries MongoDB directly — for course data, reviews, or saved plans — or triggers the image processing pipeline locally for schedule matching. Results are returned as structured JSON responses to the frontend for rendering. This fully unified architecture means the entire platform can be deployed as a single backend service alongside the React frontend, significantly reducing deployment complexity.



\section{Schedule Matcher Component}
\subsection{Description}
The Class Schedule Matcher allows EGCI students to upload screenshots of their class schedules exported from the Sky+ system and automatically identify common free time slots across a group. This removes the need for manual schedule comparison, which becomes increasingly error-prone as the number of students in a group grows. The function relies on two core image processing techniques: Optical Character Recognition (OCR) via Tesseract.js, which extracts the text-based day and time labels that define the structure of the schedule grid, and color segmentation via OpenCV.js, which detects the colored rectangular blocks representing occupied class periods. Together these two techniques convert a raw schedule screenshot into structured time slot data that can then be compared across multiple students to find common availability.
\begin{figure}[H]
	\centering
	\includegraphics[width=\textwidth]{Schedule Matcher.pdf}
	\caption{Schedule Matcher Activity Diagram}
	\label{fig:sm_ad}
\end{figure}

\subsection{Processing Pipeline}
Schedule image upload
Based on the Sky+ schedule format, the system receives screenshots that follow a consistent and predictable structure. The schedule is presented as a grid with day labels (SUN through SAT) along the vertical axis and time markers (08:00 through 20:00) along the horizontal axis. Occupied time slots are represented as solid-colored rectangular blocks placed within the grid. The grid itself sits on a white card background, providing high contrast against the colored blocks. Block colors may vary between students and between courses, but the rectangular shape and clean edges remain consistent across all inputs.

Grid Detection and Cropping
The backend first isolates the schedule grid from the rest of the screenshot. Since the Sky+ schedule card consistently appears as a white-background region with clearly defined borders, the system detects the bounding box of the grid area using edge detection in OpenCV.js and crops to that region before any further processing is done. This eliminates any interference from the background wallpaper visible around the schedule card.

Column and Row Calibration
Tesseract.js is applied to the cropped image to extract the time labels along the top of the grid (e.g. 08:00, 09:30, 11:00) and the day labels along the left side (e.g. MON, TUE). These extracted labels and their pixel positions are used to build a coordinate map that translates any horizontal pixel range into a corresponding time interval and any vertical pixel range into a corresponding day.

Color Blob Detection
OpenCV.js scans each row of the calibrated grid and identifies horizontal spans of pixels that are a colored course blob. Since the blocks are solid rectangles with clean edges against the white grid background, simple color thresholding is sufficient to detect them reliably without requiring more complex segmentation techniques. Each detected span is recorded as a start pixel and end pixel within a given day row.

Busy Map Construction and Matching 
The pixel spans from previous step are mapped to actual time slots using the coordinate map built in Step 3, producing a structured busy map for each uploaded schedule. Once all students' busy maps have been generated, the Master Busy Map algorithm performs an OR logic operation across them to identify time slots where every student is simultaneously free. If no such slot exists for the full group, the system performs a subgroup fallback by removing one student at a time and repeating the matching process, surfacing the best available options for partial groups.

\begin{figure}[H]
	\centering
	\includegraphics[width=\textwidth]{sm_pipeline.png}
	\caption{Schedule Matcher Pipeline}
	\label{fig:sm_pipeline}
\end{figure}

\subsection{Frontend}
The Schedule Matcher interface allows users to upload multiple schedule images and view the resulting common free time slots displayed in a visual timetable. Upload progress and processing status are surfaced in real time so users remain informed throughout the process.

\subsection{Backend, API, and Database}
The Express.js backend exposes an upload endpoint that receives image files from the frontend and passes them through the processing pipeline described in Section 3.2.2. Once processing is complete, the structured busy map data for each uploaded schedule is temporarily written to the schedules collection in MongoDB before the matching algorithm runs across all uploaded schedules in the session.

Each document in the schedules collection follows a lightweight structure containing a session identifier, an optional student label, a timestamp, and the extracted busy map represented as an array of occupied time slot objects. Each time slot object holds a day field and a start and end time derived from the coordinate mapping step. An example document is as follows:
\begin{figure}[H]
	\centering
	\includegraphics[width=0.8\textwidth]{sm_db}
	\caption{Schedule Matcher schedules collection example}
	\label{fig:sm_db}
\end{figure}

The session identifier groups all uploaded schedules belonging to the same matching request together, allowing the Master Busy Map algorithm to retrieve and compare all relevant documents in a single query. Once the matching results have been computed and returned to the frontend, all documents sharing that session identifier are deleted from the collection. This ensures that the schedules collection functions purely as a short-lived working space rather than a persistent data store, keeping the database lean and avoiding unnecessary retention of student schedule data.




\section{Rate-My-Course Component}
\subsection{Description}
The Rate-My-Course function allows EGCI students to rate the any available General Education courses, major required course, and engineering core course offered by MUIC and EGCI faculty. Students can submit a star rating from 0 to 5 and optionally leave a written comment, whether anonymously or not. The course catalog available for review covers Major Required Courses, Engineering Core Courses, and Computer Engineering Elective Courses offered in the MUIC 2022 catalog. Where a course has more than one section taught by different professors, the course is listed separately per professor, allowing students to give feedback that is specific to their actual classroom experience rather than the course in general. Students may also submit courses that are not present in the initial catalog, allowing the list to expand over time as new courses become available.
\begin{figure}[H]
	\centering
	\includegraphics[width=\textwidth]{Course Review.pdf}
	\caption{Rate-My-Courses Activity Diagram}
	\label{fig:rmc_ad}
\end{figure}

\subsection{Processing Pipeline}
Browe and Read Pipeline:
The pipeline begins when a student lands on the Rate-My-Course page, where the trending dashboard is loaded by querying the reviews collection for the most recently submitted reviews, grouped by course, and returning the top results to the frontend for display. If the student is looking for a specific course, they enter a course code or course name into the search bar, which sends a query to the backend courses endpoint. The backend searches the courses collection for matching documents and returns the results to the frontend as a list. The student selects a course from the results, and the course page is loaded by fetching the course metadata from the courses collection alongside a dynamic aggregation query that computes the current average rating and total review count from the reviews collection. Individual reviews for the selected course section are also retrieved and displayed in chronological order.

Review Submission pipeline:
Once the student is on a course page, they can choose to submit a review. The submission form collects a star rating, an optional written comment, and an anonymity toggle. Basic client-side validation confirms that a star rating has been selected before the form can be submitted. The review payload — containing the course code, section number, star rating, optional comment, anonymity flag, and student ID derived from the authenticated session — is then sent to the backend via a POST request to the reviews endpoint.

On the backend, the payload undergoes server-side validation to confirm that the rating is a numeric value between 0 and 5, that the referenced course and section exist in the courses collection, and that all required fields are present and correctly typed. If validation passes, the review document is written to the reviews collection with a server-generated timestamp. If any validation check fails, an error response is returned and nothing is written to the database.

Once a review is successfully stored, any subsequent request to the course page triggers a fresh aggregation query to recompute the average rating and total review count, ensuring that the displayed statistics always reflect the most current data.

\begin{figure}[H]
	\centering
	\includegraphics[width=\textwidth]{rmc_pipeline.png}
	\caption{Rate-My-Course Pipeline}
	\label{fig:rmc_pipeline}
\end{figure}

\subsection{Frontend}
The Rate-My-Course interface presents a searchable and filterable list of courses as its main view. Students can search by course code or course name using a search bar. Where a course has multiple sections under different professors, each section is displayed as a separate entry so that students can navigate directly to the professor they studied under.

The main page also features a trending courses dashboard, which highlights courses that have received the most recent review activity. This gives students an at-a-glance view of which courses are currently being discussed without requiring them to search for a specific course first, making the platform more discoverable and encouraging broader engagement across the catalog.

Each course page displays the aggregated star rating shown as both a numerical average and a visual star indicator, a breakdown of the rating distribution across all five star levels, and a chronological list of individual reviews. Reviews marked as anonymous display without a student identifier, while non-anonymous reviews show the submitter's name.

A submission form is accessible from the course page, allowing students to select a star rating by clicking on a star indicator, write an optional comment in a text field, and toggle whether their submission should appear anonymously before posting. Basic client-side validation ensures that a star rating is selected before the form can be submitted, since a rating is the only required field. The form also includes a separate flow for submitting a new course that does not yet appear in the catalog, prompting the student to provide the course code, course name, category, and professor name.

\subsection{Backend and API}
The backend exposes a set of RESTful API endpoints organized two primary resources: courses and reviews. The course endpoints handle retrieving the course page of the student’s search with its aggregated rating , and accepting new course submissions. The review endpoints handle submitting a new review for a specific course-professor pair and retrieving all reviews associated with that course section ordered by submission date and time.

All review submissions pass through server-side validation before being written to the database. The validation layer checks that the star rating is a numeric value between 0 and 5, that the course and section being reviewed exist in the database, and that the anonymity flag is a valid boolean. If any of these checks fail, the endpoint returns an appropriate error response without writing to the database. Aggregated ratings — the average score and total review count displayed on each course page — are computed dynamically from the reviews collection at read time rather than stored as fixed values, ensuring that displayed scores always reflect the most current data without requiring update operations on the course document itself. 

\subsection{Database Schema}
The Rate-My-Course function uses two database collections: courses and reviews.
The courses collection stores the metadata for each course entry. Each document contains course code, course name, category, and an array of sections objects. Each section object holds a name of the professor teaching in that section. This nested structure keeps different-professor  within the same course grouped within a single document while still allowing it to be linked to a distinct section.

The reviews section stores each individual student review submission as a separate document. Each document consists of course code, professor, course star rating, optional comment section, anonymity boolean flag, the submitter’s student ID (stored internally but masked if anonymous), and a timestamp of when the review is created.

\section{Super Planner Component}
\subsection{Description}
The Super Planner is a visual academic planning tool with interactive design that allows EGCI students to to arrange the plan for their entire academic studies by dragging and dropping course bubbles into free trimester slots. It is intended as a personal aid rather than the official enrollment system.

The list of courses available covers MUIC General Education courses, EGCI Major Required Courses, Core courses, and Elective Courses from the MUIC catalog. Each course that has prerequisite courses will have a line connecting to the prerequisite courses. There would also be an additional “suggestion” line, which connects between a course and other available courses which are suggested by the faculty to study first although the course is not connected by the prerequisite line. With the general education courses, I-design courses, and major electives courses, the website will allow users to specify what course they would take to complete the required credit of that category.
\begin{figure}[H]
	\centering
	\includegraphics[width=\textwidth]{Super Planner.pdf}
	\caption{Super Planner Activity Diagram}
	\label{fig:sp_ad}
\end{figure}

\subsection{Processing Pipeline}
The Super Planner Pipeline covers main three flows: initializing the canvas, interacting with the plan, and auto saving.

Initialization Pipeline:
When a student opens the Super Planner, the frontend simultaneously sends two requests to the backend. The first fetches the full course list from the courses collection, including each course's prerequisite relationships and faculty suggestion metadata, which is used to populate the course palette and prepare the prerequisite graph engine. The second fetches any previously saved plan from the plans collection using the authenticated student's ID.

Interaction Pipeline:
When a student drags a course icon from the palette and drops it into a semester slot, the prerequisite graph engine checks whether any of that course's prerequisites or suggested courses are already placed on the canvas and immediately renders the appropriate connecting lines. If a prerequisite course has not yet been placed, the line is rendered in a warning color to indicate an unresolved dependency.

When a student moves a course from one semester slot to another, all connecting lines associated with that course are recomputed and redrawn based on the new position, and the progress marker system reruns its check to update any warning markers that may have changed as a result of the move.

When a student marks a course as completed, the course icon updates its visual state on the canvas and the progress marker check reruns to reflect the updated academic standing.


When a student removes a course from the canvas and returns it to the palette, all connecting lines to and from that course are removed, dependent courses that relied on it as a prerequisite have their lines updated to reflect the unresolved dependency, and the progress marker system reruns to reattach a warning marker to the returned course if applicable.

Save Pipeline:
The Super Planner uses an autosave mechanism to persist the student's plan automatically without requiring a manual save action. After every interaction — a course placement, move, completion toggle, or removal — a debounce timer of approximately one to two seconds is reset. Once the student pauses and the timer expires without further changes, the canvas state is serialized from React state into a structured plan object and sent to the backend via a POST request if no plan exists yet for the student, or a PUT request if a plan already exists. The backend writes or updates the corresponding document in the plans collection and returns a success response, after which the frontend displays a brief autosave confirmation indicator such as "Saved just now" to let the student know their plan is up to date. This approach ensures that the student's work is continuously preserved without flooding the backend with requests during active interactions.

\begin{figure}[H]
	\centering
	\includegraphics[width=0.85\textwidth]{sp_pipeline.png}
	\caption{Super Planner Pipeline}
	\label{fig:sp_pipeline}
\end{figure}

\subsection{Progress Markers}
The courses that are not yet placed in the plan will also display a marker to suggest the users of their choice of course enrollment, for example, a second year course that has never been enrolled by the third year student will show a marker and a suggestion to the student that they should consider taking it as soon as possible.

\subsection{Frontend}
The Super Planner interface is built as an interactive drag-and-drop canvas within React. The canvas is organized as a grid of semester slots, each representing one trimester, with multiple semesters laid out horizontally across the screen. Course icons for Major Required courses, Engineering Core courses, and Major Elective courses are presented as a palette of floating icons around the canvas, which students can drag and drop into any semester slot they choose.

General Education courses, I-Design courses, and major elective courses are handled differently from the fixed required courses. Rather than individual course icons, these categories are represented as credit-based placeholder blocks within the plan. Students specify which courses they intend to take to fulfill the credit requirements of each category by selecting from a list within the placeholder, allowing flexibility without cluttering the canvas with every possible elective option.

When a course icon is placed into a semester slot, the interface automatically renders connecting lines to any prerequisite courses already placed on the canvas. Hard prerequisite connections are shown as solid lines, making the dependency chain immediately visible. A secondary set of dashed lines represents faculty suggestions — courses that are not formally required as prerequisites but are recommended to be completed first for better academic preparation. 

The canvas also supports marking courses as already completed, allowing students who are mid-program to accurately represent their current academic standing rather than starting from a blank slate. This allows the planner to serve both incoming students mapping out their full degree and continuing students tracking their remaining coursework.

\subsection{Backend, API, and Database}
The backend exposes API endpoints for retrieving a saved plan by student ID, creating a new plan, and updating an existing plan. An additional endpoint serves the full course list with prerequisite and suggestion metadata, which the frontend uses to initialize the course palette and render connecting lines on the canvas.

The plans collection in MongoDB stores one document per student. Each document contains the student's ID, a timestamp of the last update, and a semesters array where each element represents one trimester. Each semester object holds a semester label and a courses array containing the course codes placed in that slot along with a boolean flag indicating whether the course has been marked as completed. 

